\begin{figure}
\centering
\footnotesize
\begin{minted}[
  xleftmargin=2ex, xrightmargin=1ex, frame=lines, fontsize=\footnotesize, linenos, 
  autogobble, numbersep=1ex, framesep=1mm, breaklines]
{java}
public ReplicaHandler createTmp(...) {
  do {
    if (current == lastFoundReplica) { ... };
    // ...
    // Fault F2 (Injected)
    if ((current.genStamp() >= b.genStamp()) || !isTransfer)
      // Fault F2 (Triggered by F1 injection)
      throw new ReplicaAlreadyExistsException(...);
  } while (true);
}

class BlockReceiver {
  BlockReceiver(...){
    data.createTmp(...);
  }
}
\end{minted}
% \vspace{-3mm}
\caption{Pseudo-code demonstrating state compatibility. The fault is located inside a loop of \texttt{createTmp()}, which is invoked from \texttt{BlockReceiver()}. Code simplified from \texttt{FsDatasetImpl} in HDFS.}
% \vspace{-2mm}
\label{fig:state-compatibility}
\end{figure}
